\section{Experimental Protocol}
\label{sec:experiments}

\subsection{Evidence Design and Datasets}

The development archive contains 15 valid datasets, 17 complete configurations,
and three seeds. The main comparison reports 16 primary configurations and
excludes the exploratory scVICAR-T contrastive extension, which is retained in
the supplement. Each reported configuration completed 45 runs using its
method-specific protocol. The final scVICAR-F
and scVICAR-T configurations were selected during this development stage. The
benchmark measures performance across the broad scMAE dataset collection. The
confirmatory cohort contains six additional datasets selected before the fixed
and topology-adaptive models were evaluated (Table~\ref{tab:data}).
We remove labels whose normalized value is \texttt{unknown}, \texttt{nan}, or
\texttt{unassigned}, followed by classes with fewer than ten cells.  The
unfiltered-label analysis is retained as sensitivity analysis.  The
\texttt{hrvatin\_geo} artifact is excluded because its labels are invalid for
formal aggregation and its cells overlap with \texttt{hrvatin}.

The 15 development datasets are Bach, Baron, Guo, Limb Muscle, Macosko,
Melanoma 5K, Pollen, Quake 10x Spleen, Quake Smart-seq2 Lung, Shekhar,
Tosches, Wang, Young, Hrvatin, and worm neuron cell. They span pancreas,
retina, melanoma, lung, spleen, muscle, embryonic, and neuronal systems and
include plate-based and droplet sequencing protocols. Dataset provenance,
dimensions, label fields, and preprocessing are listed in the supplement.

\begin{table}[t]
\centering
\caption{Frozen confirmatory datasets after primary filtering.}
\label{tab:data}
\small
\begin{tabular}{lrrr}
\toprule
Dataset & Cells & Genes & $K$\\
\midrule
Blood--Bone Marrow & 15,486 & 61,497 & 30\\
Human Pancreas 1 & 2,282 & 61,497 & 6\\
Human Pancreas 3 & 8,562 & 20,125 & 13\\
Mouse Pancreas 1 & 1,865 & 14,878 & 10\\
PRJNA895163 & 50,380 & 25,415 & 12\\
Tabula Sapiens Pancreas & 14,116 & 61,497 & 16\\
\bottomrule
\end{tabular}
\end{table}

The dataset manifest records source identifiers, label fields, filtering rules,
dimensions, and input checksums. Blood--Bone Marrow derives from the hematopoietic
reference map~\cite{triana2021hematopoietic}; Human Pancreas 1 is from Enge et
al.~\cite{enge2017pancreas}; PRJNA895163 is the cassava tuberous-root atlas of
Song et al.~\cite{song2022cassava}; and the Tabula Sapiens subset derives from
the multi-organ atlas~\cite{tabulasapiens2022atlas}.  Baron human/mouse pancreas
citations and processing-stage counts are detailed in
Section~\ref{sec:downstream} and the supplement.
Per-cell labels are reserved for post-training evaluation and prespecified
diagnostic interventions. Clean-view optimization, graph construction,
checkpoint selection, and run scheduling use expression data and frozen
configuration fields.

\subsection{Matched Variants and Baselines}

We compare six variants under the common backbone in Section~\ref{sec:method}:
NoMix, random mixing, scVICAR-F, topology edge-only, topology gate-only, and
full scVICAR-T.  NoMix is the matched scMAE objective; random mixing tests
whether arbitrary convex corruption explains any gain.  Component variants
separate neighbor affinity from the cell-specific mixing coefficient.

External references comprise PCA+KMeans, the original scMAE implementation,
scVI, scDeepCluster, scDCC, scRCL, and graph/deep-clustering methods supported by
the benchmark adapters~\cite{lopez2018scvi,tian2019scdeepcluster,tian2021scdcc,peng2024scrcl}.
The 15-dataset development table additionally includes Louvain, Leiden, SC3,
DEC, scNAME, scziDesk, scCDCG, and PhytoCluster
~\cite{blondel2008louvain,traag2019leiden,kiselev2017sc3,xie2016dec,wan2022scname,chen2020sczidesk,xu2024sccdcg,wang2025phytocluster}.
Historical values are reused only when the input and evaluation protocol match;
otherwise the baseline is rerun. Method-family conclusions are based on the
matched variants, avoiding attribution of backbone changes to vicinal recovery. The
published architectures of scDCC, scDeepCluster, and scRCL require $K$ during
clustering-oriented training; we therefore mark them explicitly as known-$K$
training baselines.  Per-cell reference labels are withheld from their gradient
updates and checkpoint selection.

\subsection{Evaluation and Frozen Hyperparameters}

Model seeds are 42, 2024, and 3407.  The six datasets, six matched variants,
and three seeds give \FormalRunCount{} primary runs.  The primary endpoint is
ARI from known-$K$ KMeans with \texttt{n\_init=20}.  NMI, assignment accuracy,
macro-F1, silhouette, runtime, and peak GPU memory are secondary.  Fixed
resolution 1.0 Leiden is reported separately with a fixed neighborhood graph.
The six matched scVICAR variants use known $K$ only during post-hoc evaluation.
The three external known-$K$ baselines use it during their published
clustering-oriented training procedure.

The graph uses $k=5$, four sampled neighbors, $\tau=0.2$, and a maximum gate of
0.1.  All models train for 80 epochs with batch size 256, learning rate
$10^{-3}$, mask ratio 0.4, and vicinal weight 0.3 where applicable. One
configuration is used for all six datasets. The primary model sets the
contrastive weight to zero and contains a static graph, analytic coefficient,
and the losses defined in Eq.~\eqref{eq:objective}.

\subsection{Graph-Contamination and Estimator Tests}

On Blood--Bone Marrow, Human Pancreas 3, and PRJNA895163, we replace 0, 25, 50,
75, or 100\% of each graph row with deliberately cross-label neighbors. Labels
construct this isolated stress intervention. We compare scVICAR-F with scVICAR-T by absolute performance
and degradation from 0\% contamination.  A separate aggregation ablation
compares the implemented sampling-then-reweighting estimator in
Eq.~\eqref{eq:realized-kernel}, an unbiased uniform average of sampled
neighbors, and full-neighborhood aggregation.

After training, we compute Spearman association and AUROC between affinity or
mixing scores and same-type edge purity. We also relate the coefficient to
class frequency and per-class recall. These calculations use the frozen graph,
model, and checkpoint.

\subsection{Statistical Analysis}

The dataset is the independent unit. For downstream tasks, split seeds are first
averaged within model run; model seeds are then averaged within dataset.  The
three prespecified contrasts are F--NoMix, T--NoMix, and T--F.  We report paired
dataset-level mean and median effects, win/tie/loss counts, hierarchical
bootstrap 95\% confidence intervals, paired sign-flip permutation tests, and
Holm-adjusted $p$-values. Cells, splits, and seeds are nested computational
observations. Every seed is retained.

\subsection{Reproducibility and Traceability}

Every reported number and figure resolves to a frozen run and configuration.
The supplement describes checksums, run identifiers, completion criteria,
software identities, and storage provenance.
